% #############################################################################
% Abstract
% !TEX root = ../../Main.tex
% #############################################################################
\acresetall
\noindent Reinforcement learning for currency trading is usually tested on one pair, on bar data and without real costs. This thesis trains one Deep Deterministic Policy Gradient (DDPG) agent, with one design and no tuning per pair, on each of the seven major currency pairs, over eleven years of tick data from 2015 to 2025. The agent is implemented as in the original DDPG paper: an off-policy actor-critic with actor and critic networks, slowly updated target copies of both, a replay buffer and Ornstein-Uhlenbeck exploration noise, with the published hyperparameters where they apply. The actor maps 38 dimensionless features from sixteen technical indicators to one continuous action in $[-1, 1]$, whose sign sets the direction of the position and whose magnitude sets its size within a volatility-based risk budget. The critic learns the value of that action, and the reward is the log return of account equity. Four documented changes adapt it to this problem: smaller networks with layer normalisation, a penalty against output saturation, a scaled reward and mirrored training episodes. The agent is trained and evaluated in a tick-level backtesting engine, built for this work to reproduce cTrader's backtester, that replays 1.6 billion quotes and charges the real costs of a raw-spread, swap-free retail account: the spread quoted at each fill, a commission of 3.5 US dollars per lot per side and a swap of zero. Each pair is trained with an anchored walk-forward, the standard validation method for trading systems: seven folds validate on 2018 to 2024, 2025 is held out as the test year, and a model is kept only if it passes eleven checks on return and behaviour. All seven models are profitable net of all costs. They returned between \ResultStitchedMin{} and \ResultStitchedMax{} over the walk-forward years, between \ResultTestMin{} and \ResultTestMax{} in the held-out year and between \ResultFullMin{} and \ResultFullMax{} over the full period, with Sharpe ratios of \ResultSharpeMin{} to \ResultSharpeMax{} and a yearly alpha of \ResultAlphaMin{} to \ResultAlphaMax{}, and each beats its own pair in all three windows. As an equal-weight portfolio they returned \ResultPortfolioReturn{} against \ResultBenchmarkReturn{} for the US 500 index, with a maximum drawdown of \ResultPortfolioDrawdown{} against \ResultBenchmarkDrawdown{} and a Sharpe ratio of \ResultPortfolioSharpe{} against \ResultBenchmarkSharpe{}, because the seven models make and lose money at different times.
